% Three compact state-selection diagrams paired with the intervention matrix.
\begin{tikzpicture}[x=1cm,y=1cm,font=\footnotesize, % nocheck: schematic layout
 state/.style={circle,draw=black!40,fill=white,minimum size=4mm,inner sep=0pt}, % nocheck: schematic layout
 selected/.style={state,draw=palB,fill=palB!18,line width=.9pt}, % nocheck: schematic layout
 flow/.style={-{Stealth[length=3pt]},draw=black!50}, % nocheck: schematic layout
 head/.style={anchor=west,font=\footnotesize\bfseries},
 note/.style={anchor=west,align=left,font=\footnotesize}]
\useasboundingbox (0,.3) rectangle (7.5,-5.35); % nocheck: schematic layout
\fill[palB!5,rounded corners=2pt] (0,.3) rectangle (7.5,-1.35); % nocheck: schematic layout
\node[head] at (.18,.02) {First divergence}; % nocheck: schematic layout
\node[anchor=east,text=black!60] at (7.3,.02) {\S\ref{sec:value}, \S\ref{sec:sameprec}}; % nocheck: schematic layout
\node[state] (a1) at (.5,-.55) {}; % nocheck: schematic layout
\node[state] (a2) at (1.3,-.55) {}; % nocheck: schematic layout
\node[selected] (a3) at (2.1,-.55) {}; % nocheck: schematic layout
\draw[flow] (a1)--(a2); \draw[flow] (a2)--(a3); % nocheck: schematic layout
\node[note] at (2.65,-.55) {Stop when B and Q first differ}; % nocheck: schematic layout
\node[note,text=black!65] at (.18,-1.06) {Q's history; compare tool structure, not wording.}; % nocheck: schematic layout
\fill[black!3,rounded corners=2pt] (0,-1.55) rectangle (7.5,-3.2); % nocheck: schematic layout
\node[head] at (.18,-1.85) {Ordinary sampling}; % nocheck: schematic layout
\node[anchor=east,text=black!60] at (7.3,-1.85) {\S\ref{sec:trigger}}; % nocheck: schematic layout
\node[state] (b1) at (.5,-2.4) {}; % nocheck: schematic layout
\node[selected] (b2) at (1.3,-2.4) {}; % nocheck: schematic layout
\node[state] (b3) at (2.1,-2.4) {}; % nocheck: schematic layout
\draw[flow] (b1)--(b2); \draw[flow] (b2)--(b3); % nocheck: schematic layout
\node[note] at (2.65,-2.4) {Select by a fixed hash}; % nocheck: schematic layout
\node[note,text=black!65] at (.18,-2.91) {Selection uses neither disagreement nor outcome.}; % nocheck: schematic layout
\fill[palQP!7,rounded corners=2pt] (0,-3.4) rectangle (7.5,-5.35); % nocheck: schematic layout
\node[head] at (.18,-3.7) {Common-state comparison}; % nocheck: schematic layout
\node[anchor=east,text=black!60] at (7.3,-3.7) {\S\ref{sec:sameprec}}; % nocheck: schematic layout
\node[draw=black!30,rounded corners=2pt,align=center,inner sep=4pt] (h) at (1.25,-4.37) {One sampled\\history}; % nocheck: schematic layout
\node[text=palB,font=\footnotesize\bfseries] (c1) at (3.15,-4.05) {B}; % nocheck: schematic layout
\node[text=palQ,font=\footnotesize\bfseries] (c2) at (3.15,-4.42) {Q}; % nocheck: schematic layout
\node[text=palQP,font=\footnotesize\bfseries] (c3) at (3.15,-4.79) {Q$'$}; % nocheck: schematic layout
\draw[flow] (h.east)--(c1.west);\draw[flow] (h.east)--(c2.west);\draw[flow] (h.east)--(c3.west); % nocheck: schematic layout
\node[note,text width=3.2cm] at (3.85,-4.4) {Compare actions under\\the same continuation}; % nocheck: schematic layout
\node[note,text=black!65] at (.18,-5.1) {Does BF16 add value over another W4 checkpoint?}; % nocheck: schematic layout
\end{tikzpicture}
